% ----------------------------------------------------------------------
%  3.4.3. Атомная теория диффузии в металлах
% ----------------------------------------------------------------------
\subsection{Атомная теория диффузии в металлах}\label{ssec:diff-atomic}

Континуальный подход даёт общие соотношения, но ничего не говорит о
том, \emph{как} движутся атомы. Атомистический подход связывает
макроскопические величины (потоки, коэффициенты диффузии) с атомными
характеристиками (частотами и длинами скачков, энергиями дефектов).

\paragraph{Механизмы диффузии}
Атомы в кристалле колеблются около положений равновесия; при нагреве
колебания становятся настолько сильными, что атом время от времени
\emph{перескакивает} в соседнее положение. Основное предположение
теории: диффузия~— последовательность таких скачков между положениями
равновесия. По типу скачка различают несколько механизмов
(рис.~\ref{fig:diff-mech}).

\newcommand{\DiffLat}{\foreach \x in {0,...,4}{\foreach \y in {0,...,3}{\node[atom=2.5mm] at (\x,\y) {};}}}
\begin{figure}[ht]
\centering
\begin{tikzpicture}[scale=0.5,
  hop/.style={-{Stealth[length=1.7mm]}, line width=0.8pt, fmRed}]
  % а — прямой обмен
  \begin{scope}
    \DiffLat \node[atomB=2.5mm] at (2,2) {}; \node[atomB=2.5mm] at (2,1) {};
    \draw[hop] (2.25,2) to[bend left=50] (2.25,1);
    \draw[hop] (1.75,1) to[bend left=50] (1.75,2);
    \node[font=\small] at (2,-1.1) {\textit{а}};
  \end{scope}
  % б — кольцевой
  \begin{scope}[shift={(6.3,0)}]
    \DiffLat
    \foreach \p in {(1,1),(2,1),(2,2),(1,2)}{\node[atomB=2.5mm] at \p {};}
    \draw[hop] (1.2,1.15) -- (1.8,1.15); \draw[hop] (2.15,1.2) -- (2.15,1.8);
    \draw[hop] (1.8,1.85) -- (1.2,1.85); \draw[hop] (0.85,1.8) -- (0.85,1.2);
    \node[font=\small] at (2,-1.1) {\textit{б}};
  \end{scope}
  % в — прямой межузельный
  \begin{scope}[shift={(12.6,0)}]
    \DiffLat
    \node[atomI=1.5mm] at (1.5,1.5) {};
    \draw[hop] (1.65,1.5) -- (2.35,1.5);
    \node[atomI=1.5mm, fill opacity=0.35, draw opacity=0.35] at (2.5,1.5) {};
    \node[font=\small] at (2,-1.1) {\textit{в}};
  \end{scope}
  % г — непрямой (вытеснение)
  \begin{scope}[shift={(0,-5.6)}]
    \DiffLat
    \node[atomI=2.5mm] at (1.5,1.5) {};
    \draw[hop] (1.6,1.4) -- (1.9,1.1);
    \draw[hop] (2.15,0.95) -- (2.4,0.6);
    \node[atomB=2.5mm] at (2,1) {};
    \node[font=\small] at (2,-1.1) {\textit{г}};
  \end{scope}
  % д — краудион
  \begin{scope}[shift={(6.3,-5.6)}]
    \foreach \x in {0,...,4}{\foreach \y in {0,1,3}{\node[atom=2.5mm] at (\x,\y) {};}}
    \foreach \i in {0,...,5}{\node[atomB=2.5mm] at ({\i*0.8},2) {};}
    \draw[hop] (0,2.55) -- (3.6,2.55);
    \node[font=\small] at (2,-1.1) {\textit{д}};
  \end{scope}
  % е — вакансионный
  \begin{scope}[shift={(12.6,-5.6)}]
    \foreach \x in {0,...,4}{\foreach \y in {0,...,3}{
      \ifnum\x=2 \ifnum\y=1 \node[vac=2.5mm] at (\x,\y) {};\else\node[atom=2.5mm] at (\x,\y) {};\fi
      \else\node[atom=2.5mm] at (\x,\y) {};\fi}}
    \node[atomB=2.5mm] at (1,1) {};
    \draw[hop] (1.25,1.2) to[bend left=30] (1.8,1.2);
    \node[font=\small] at (2,-1.1) {\textit{е}};
  \end{scope}
\end{tikzpicture}
\caption{Механизмы диффузии: \textit{а}~— прямой обмен местами;
\textit{б}~— кольцевой; \textit{в}~— прямой межузельный (малый атом
перескакивает между междоузлиями); \textit{г}~— непрямой межузельный
(межузельный атом вытесняет атом из узла в соседнее междоузлие);
\textit{д}~— краудион (лишний атом в плотноупакованном ряду);
\textit{е}~— вакансионный}
\label{fig:diff-mech}
\end{figure}

\begin{itemize}
\item \term{Обменные механизмы.} При \emph{прямом обмене} два соседних
атома меняются местами. В плотных упаковках это маловероятно: нужны
огромные искажения решётки. При \emph{кольцевом} механизме
одновременно смещаются несколько атомов, искажения меньше, но нужно
согласованное движение многих атомов~— тоже маловероятно.
\item \term{Межузельные механизмы.} \emph{Прямой}: атом перескакивает
из междоузлия в междоузлие~— так диффундируют малые атомы примесей (H, C,
N). \emph{Непрямой}: межузельный атом, близкий по размеру к атомам
решётки, вытесняет атом из узла в соседнее междоузлие и занимает его
место; движется не атом, а межузельная конфигурация. В металлах
собственный межузельный атом обычно образует \emph{гантель}~— два
атома, симметрично расположенные относительно узла (вдоль
$\langle100\rangle$ в ГЦК и $\langle110\rangle$ в ОЦК). Такие механизмы
важны при облучении, когда рождаются пары Френкеля.
\emph{Краудион}~— лишний атом, «втиснутый» в плотноупакованный ряд;
возникает чаще всего при облучении и низких температурах, при нагреве
превращается в гантель.
\item \term{Вакансионный механизм}~— главный для самодиффузии и
растворов замещения. Вблизи точки плавления доля вакантных узлов
достигает $10^{-4}$--$10^{-3}$ (п.~2.1). Атом перескакивает в соседнюю
вакансию (NN-скачок, от англ.\ \emph{nearest neighbour}), реже~— в
следующую за соседней (NNN). Вклад могут давать и дивакансии (при
высоких температурах), и комплексы «растворённый атом + вакансия».
\end{itemize}

\paragraph{Откуда берётся энергия активации}
Скорость диффузии при любом «дефектном» механизме пропорциональна двум
величинам.
\begin{shagi}
\item \emph{Вероятность найти дефект} рядом с атомом (больцмановский
множитель, п.~2.1):
$n=\exp(-G_f/kT)=n_0\exp(-E_f/kT)$, $n_0=\exp(S_f/k)$, где $E_f$ и $S_f$~—
энергия и энтропия образования дефекта.
\item \emph{Частота скачка} через барьер (рис.~\ref{fig:barrier}): атом
должен пройти через \emph{седловую точку} с повышенной энергией,
$w=w_0\exp(-E_m/kT)$, $w_0=w_0'\exp(S_m/k)$\ispr{В оригинале (с.~444)
энтропийный множитель записан как $\exp(S_m/kT)$; энтропия входит в
виде $\exp(S_m/k)$.}, где $E_m$ и $S_m$~— энергия и энтропия миграции.
\item Частота скачков атома в данное соседнее положение $\Gamma=wn$; если
таких положений $z$, общая частота
\end{shagi}
\begin{keyeq}\label{eq:jump-freq}
\nu=z\Gamma=\nu_0\exp\!\left(-\frac{Q}{kT}\right),\qquad Q=E_f+E_m .
\end{keyeq}
$Q$~— \term{энергия активации диффузии}. Для меди, например,
$E_f\approx1{,}3$~эВ, $E_m\approx0{,}7$~эВ, и $Q\approx2{,}0$~эВ~— это
согласуется с опытом. Для малых атомов внедрения, которые всегда сидят
в междоузлиях, $n=1$ и $E_f=0$: энергия активации равна одной лишь
энергии миграции, поэтому углерод в железе диффундирует несравненно
быстрее самого железа. Кольцевой механизм дефектов не требует, и для
него тоже $E_f=0$.

\begin{figure}[ht]
\centering
\begin{tikzpicture}
\begin{axis}[width=0.6\textwidth, height=4.4cm, xmin=-0.5, xmax=2.5, ymin=-0.15, ymax=1.35,
  xlabel={положение атома}, ylabel={энергия}, xtick={0,1,2}, xticklabels={узел,узел,узел},
  ytick=\empty, axis lines=left]
  \addplot[fmBlue, thick, domain=-0.5:2.5, samples=150] {0.5*(1-cos(deg(2*pi*x)))};
  \draw[dimline] (axis cs:0.5,0) -- (axis cs:0.5,1) node[midway, right, font=\small]{$E_m$};
  \node[atomB=3mm] at (axis cs:0,0.09) {};
  \draw[-{Stealth[length=2mm]}, fmRed] (axis cs:0.05,0.25) to[bend left=40] (axis cs:0.95,0.25);
  \node[font=\scriptsize, above] at (axis cs:0.5,1.03) {седловая точка};
\end{axis}
\end{tikzpicture}
\caption{Энергия атома при перескоке в соседнее положение: чтобы
перейти, атом должен преодолеть барьер высотой $E_m$}
\label{fig:barrier}
\end{figure}

\paragraph{Случайные скачки дают закон Фика}
Рассмотрим две соседние атомные плоскости $\Pi_1$ и $\Pi_2$,
перпендикулярные оси $x$, на расстоянии $\lambda$ друг от друга
(рис.~\ref{fig:planes}). На единице площади плоскостей $n_1$ и $n_2$
атомов диффузанта; частота скачка в соседнюю плоскость~— $\Gamma_0$
(одинакова в обе стороны).
\begin{shagi}
\item Результирующий поток: $J=n_1\Gamma_0-n_2\Gamma_0=(n_1-n_2)\Gamma_0$.
\item Поверхностная плотность связана с объёмной концентрацией: $n=\lambda C$,
поэтому $n_1-n_2=-\lambda^2\,\partial C/\partial x$.
\item Получаем $J=-\lambda^2\Gamma_0\,\partial C/\partial x$~— первый закон
Фика.
\end{shagi}
Обозначив $\Gamma_0=\gamma\nu$ (множитель $\gamma$ учитывает геометрию
решётки и механизм; в простой кубической решётке $\gamma=1/6$),
получаем коэффициент диффузии и его температурную зависимость~—
\term{уравнение Аррениуса}:
\begin{keyeq}\label{eq:arrhenius}
D^*=\gamma\lambda^2\nu=D_0\exp\!\left(-\frac{Q}{kT}\right),\qquad D_0=\gamma\lambda^2\nu_0 .
\end{keyeq}
Для вакансионного механизма, например,
$D^*=\gamma\lambda^2\nu_0'\exp\frac{S_f+S_m}{k}\exp\left(-\frac{H_f+H_m}{kT}\right)$,
для межузельного~— $D^*=\gamma\lambda^2\nu_0'\exp\frac{S_m}{k}\exp\left(-\frac{H_m}{kT}\right)$.

\begin{figure}[ht]
\centering
\begin{subfigure}[b]{0.42\textwidth}\centering
\begin{tikzpicture}[scale=0.62]
  \foreach \x/\l in {0/{$\Pi_1$}, 2.2/{$\Pi_2$}}{
    \draw[fmGray, line width=0.6pt] (\x,-0.4) -- (\x,4.4);
    \node[font=\small, above] at (\x,4.4) {\l};}
  \foreach \y in {0,0.8,...,4}{\node[atomB=3mm] at (0,\y) {};}
  \foreach \y in {0.4,2.0,3.6}{\node[atomB=3mm] at (2.2,\y) {};}
  \draw[-{Stealth[length=2mm]}, fmRed, thick] (0.3,2.6) -- (1.9,2.6) node[midway, above, font=\scriptsize]{$n_1\Gamma_0$};
  \draw[-{Stealth[length=2mm]}, fmBlue, thick] (1.9,1.3) -- (0.3,1.3) node[midway, below, font=\scriptsize]{$n_2\Gamma_0$};
  \draw[dimline] (0,-0.9) -- node[below, font=\small]{$\lambda$} (2.2,-0.9);
  \draw[-{Stealth[length=2mm]}] (3.2,2) -- (4.4,2) node[right, font=\small]{$x$};
\end{tikzpicture}
\caption{}
\end{subfigure}\hfill
\begin{subfigure}[b]{0.5\textwidth}\centering
\begin{tikzpicture}[scale=0.21, walk/.style={fmBlue, line width=0.5pt, line join=round}]
  % сгенерировано tools/gen_diffusion.py — не править вручную
\draw[walk] (0,0) -- (0,1) -- (-1,1) -- (-1,0) -- (0,0) -- (1,0) -- (2,0) -- (2,1) -- (3,1) -- (2,1) -- (3,1) -- (4,1) -- (4,0) -- (4,-1) -- (5,-1) -- (4,-1) -- (5,-1) -- (5,-2) -- (6,-2) -- (7,-2) -- (6,-2) -- (7,-2) -- (7,-3) -- (8,-3) -- (7,-3) -- (8,-3) -- (7,-3) -- (7,-2) -- (7,-3) -- (6,-3) -- (7,-3) -- (7,-2) -- (6,-2) -- (7,-2) -- (6,-2) -- (6,-1) -- (7,-1) -- (8,-1) -- (9,-1) -- (8,-1) -- (8,-2) -- (8,-3) -- (8,-2) -- (8,-3) -- (8,-4) -- (8,-3) -- (8,-2) -- (7,-2) -- (6,-2) -- (5,-2) -- (6,-2) -- (6,-1) -- (6,-2) -- (6,-1) -- (6,-2) -- (6,-1) -- (7,-1) -- (8,-1) -- (8,-2) -- (7,-2) -- (7,-1) -- (6,-1) -- (6,-2) -- (6,-3) -- (7,-3) -- (8,-3) -- (8,-2) -- (8,-1) -- (8,0) -- (8,-1) -- (8,-2) -- (9,-2) -- (10,-2) -- (10,-1) -- (10,-2) -- (11,-2) -- (12,-2) -- (12,-1) -- (12,-2) -- (12,-1) -- (12,-2) -- (12,-1) -- (13,-1) -- (13,-2) -- (13,-1) -- (12,-1) -- (13,-1) -- (13,-2) -- (14,-2) -- (13,-2) -- (13,-1) -- (12,-1) -- (11,-1) -- (11,-2) -- (11,-3) -- (11,-4) -- (12,-4) -- (11,-4) -- (11,-5) -- (11,-6) -- (11,-5) -- (10,-5) -- (10,-6) -- (10,-5) -- (10,-6) -- (10,-5) -- (10,-6) -- (9,-6) -- (8,-6) -- (9,-6) -- (8,-6) -- (7,-6) -- (6,-6) -- (5,-6) -- (6,-6) -- (6,-7) -- (5,-7) -- (5,-6) -- (5,-5) -- (6,-5) -- (5,-5) -- (5,-6) -- (5,-5) -- (5,-4) -- (4,-4) -- (5,-4) -- (5,-5) -- (5,-6) -- (5,-7) -- (5,-8) -- (5,-9) -- (6,-9) -- (6,-10) -- (6,-11) -- (7,-11) -- (6,-11) -- (7,-11) -- (6,-11) -- (6,-12) -- (5,-12) -- (6,-12) -- (6,-11) -- (7,-11) -- (8,-11) -- (9,-11) -- (8,-11) -- (9,-11) -- (9,-10) -- (10,-10) -- (11,-10) -- (10,-10) -- (10,-11) -- (9,-11) -- (9,-10) -- (9,-9) -- (9,-8) -- (9,-9) -- (10,-9) -- (11,-9) -- (11,-10) -- (11,-11) -- (11,-12) -- (11,-13) -- (11,-12) -- (12,-12) -- (11,-12) -- (12,-12) -- (12,-11) -- (12,-10) -- (12,-11) -- (11,-11) -- (12,-11) -- (11,-11) -- (11,-10) -- (10,-10) -- (11,-10) -- (11,-9) -- (12,-9) -- (12,-8) -- (12,-7) -- (11,-7) -- (11,-6) -- (10,-6) -- (10,-5) -- (9,-5) -- (8,-5) -- (7,-5) -- (7,-6) -- (6,-6) -- (5,-6) -- (5,-7) -- (5,-6) -- (6,-6) -- (7,-6) -- (7,-5) -- (7,-6) -- (7,-5) -- (6,-5) -- (6,-4) -- (6,-5) -- (6,-4) -- (6,-3) -- (7,-3) -- (6,-3) -- (7,-3) -- (6,-3) -- (6,-4) -- (5,-4) -- (5,-3) -- (4,-3) -- (4,-4) -- (5,-4) -- (5,-5) -- (5,-4) -- (6,-4) -- (7,-4) -- (7,-5) -- (6,-5) -- (6,-6) -- (5,-6) -- (5,-7) -- (5,-6) -- (6,-6) -- (6,-7) -- (6,-8) -- (6,-9) -- (7,-9) -- (6,-9) -- (5,-9) -- (4,-9) -- (5,-9) -- (4,-9) -- (4,-10) -- (3,-10) -- (3,-11) -- (3,-10) -- (2,-10) -- (1,-10) -- (2,-10) -- (3,-10) -- (4,-10) -- (3,-10) -- (3,-11) -- (2,-11) -- (1,-11) -- (2,-11) -- (2,-10) -- (1,-10) -- (1,-9) -- (0,-9) -- (0,-8) -- (0,-7) -- (0,-8) -- (-1,-8) -- (0,-8) -- (0,-7) -- (0,-8) -- (0,-9) -- (-1,-9) -- (-2,-9) -- (-1,-9) -- (-1,-10) -- (-2,-10) -- (-1,-10) -- (-2,-10) -- (-3,-10) -- (-4,-10) -- (-4,-11) -- (-3,-11) -- (-2,-11) -- (-2,-10) -- (-2,-11) -- (-1,-11) -- (0,-11) -- (-1,-11) -- (-2,-11) -- (-2,-10) -- (-1,-10) -- (0,-10) -- (0,-11) -- (1,-11) -- (2,-11) -- (2,-12) -- (2,-11) -- (1,-11) -- (1,-10) -- (1,-11) -- (1,-12) -- (0,-12) -- (0,-11) -- (-1,-11) -- (-1,-12) -- (-2,-12) -- (-2,-13) -- (-1,-13) -- (-1,-14) -- (-1,-15) -- (-1,-14) -- (0,-14) -- (-1,-14) -- (-1,-15) -- (0,-15) -- (-1,-15) -- (-1,-14) -- (0,-14) -- (-1,-14) -- (-1,-13) -- (-2,-13) -- (-2,-12) -- (-3,-12) -- (-3,-13) -- (-4,-13) -- (-3,-13) -- (-3,-14) -- (-3,-15) -- (-4,-15) -- (-5,-15) -- (-6,-15) -- (-6,-16) -- (-6,-17) -- (-6,-16) -- (-6,-17) -- (-7,-17) -- (-7,-16) -- (-7,-15) -- (-6,-15) -- (-6,-14) -- (-5,-14) -- (-5,-13) -- (-5,-14) -- (-5,-15) -- (-4,-15) -- (-4,-16) -- (-4,-15) -- (-4,-14) -- (-3,-14) -- (-2,-14) -- (-3,-14) -- (-2,-14) -- (-1,-14) -- (-1,-13) -- (-1,-12) -- (0,-12) -- (-1,-12) -- (-1,-11) -- (-2,-11) -- (-2,-12) -- (-2,-11) -- (-2,-12) -- (-3,-12) -- (-3,-13) -- (-3,-12) -- (-2,-12) -- (-2,-11) -- (-1,-11) -- (-2,-11) -- (-2,-12) -- (-1,-12) -- (-1,-11) -- (0,-11) -- (1,-11) -- (1,-10) -- (2,-10) -- (1,-10) -- (2,-10) -- (2,-9) -- (3,-9) -- (3,-10) -- (4,-10) -- (4,-9) -- (4,-10) -- (4,-9) -- (3,-9) -- (4,-9) -- (3,-9) -- (4,-9) -- (3,-9) -- (3,-8) -- (4,-8) -- (3,-8) -- (2,-8) -- (2,-7) -- (2,-6) -- (1,-6) -- (1,-5) -- (1,-6) -- (0,-6) -- (0,-5) -- (0,-4) -- (1,-4) -- (1,-3) -- (2,-3) -- (3,-3) -- (4,-3) -- (3,-3) -- (3,-4) -- (2,-4) -- (2,-5) -- (3,-5) -- (3,-6) -- (3,-7);
\coordinate (walkEnd) at (3, -7);

  \fill[fmRed] (0,0) circle (0.45);
  \fill[fmGreen!60!black] (walkEnd) circle (0.45);
  \draw[-{Stealth[length=2mm]}, fmRed, thick] (0,0) -- (walkEnd);
\end{tikzpicture}
\caption{}
\end{subfigure}
\caption{\textit{а}~— скачки между соседними атомными плоскостями:
поток в одну сторону пропорционален $n_1$, в другую~— $n_2$;
\textit{б}~— случайное блуждание (расчёт, 400 скачков): атом проходит
путь в 400 межатомных расстояний, а удаляется от начала (красная
стрелка) лишь на несколько расстояний~— в среднем на $\sqrt{400}=20$}
\label{fig:planes}
\end{figure}

\paragraph{Соотношение Эйнштейна}
Ту же связь можно получить, следя за одним атомом. Если скачки
хаотичны, средний квадрат смещения атома вдоль $x$ растёт
пропорционально времени:
\begin{keyeq}\label{eq:einstein-mob}
\langle X^2\rangle=2D^*t\qquad\Bigl(\text{в трёх измерениях}\ \langle R^2\rangle=6D^*t\Bigr).
\end{keyeq}
\begin{podrobnee}[вывод]
Пусть $W(X,\tau)$~— вероятность того, что за время $\tau$ атом сместится
на $X$. Концентрация в момент $t+\tau$:
$C(x,t+\tau)=\sum_XC(x-X,t)W(X,\tau)$. Разложим обе части в ряд при
малых $\tau$ и $X$: $C+\tau\frac{\partial C}{\partial t}=
\sum_X\left[C-X\frac{\partial C}{\partial x}+\frac{X^2}2\frac{\partial^2C}{\partial x^2}\right]W$.
Учитывая $\sum_XW=1$ и $\langle X\rangle=0$ для случайных блужданий,
получаем $\frac{\partial C}{\partial t}=\frac{\langle X^2\rangle}{2\tau}\frac{\partial^2C}{\partial x^2}$~—
второй закон Фика с $D^*=\lim\limits_{\tau\to0}\langle X^2\rangle/2\tau$.
\end{podrobnee}

\begin{primer}[Сколько прыгает атом меди]
Самодиффузия меди: $D_0=2\cdot10^{-5}$~м$^2$/с, $Q=197$~кДж/моль. При
1300~К (чуть ниже $\Tpl=1356$~К)
$D^*=2\cdot10^{-5}\exp(-197\,000/8{,}31\cdot1300)\approx2{,}4\cdot10^{-13}$~м$^2$/с.
\begin{shagi}
\item В ГЦК-решётке $D^*=\frac16\nu\lambda^2$ с $\lambda=0{,}256$~нм, откуда частота
скачков $\nu=6D^*/\lambda^2\approx2\cdot10^7$~с$^{-1}$.
\item За час атом совершает $\approx8\cdot10^{10}$ скачков и проходит путь
$8\cdot10^{10}\cdot0{,}256~\text{нм}\approx20$~м.
\item А удаляется от исходного места лишь на $\sqrt{6D^*t}\approx70$~мкм
(здесь $6D^*t=6\cdot2{,}4\cdot10^{-13}\cdot3600\approx5\cdot10^{-9}$~м$^2$).
\end{shagi}
При комнатной температуре $D^*\sim10^{-40}$~м$^2$/с: атом прыгает
примерно раз в несколько миллионов лет.
\end{primer}

\paragraph{Корреляционные эффекты}
Мы считали, что направление каждого скачка не зависит от предыдущих.
При вакансионном механизме это не так. Поменявшись местами с вакансией,
атом оказывается рядом с ней, и вероятность того, что следующий его
скачок будет \emph{обратно} в ту же вакансию, повышена. Такие пары
скачков «вперёд--назад» ничего не дают для диффузии. Поэтому
эффективная частота скачков уменьшается на \term{корреляционный
множитель} $f$, и
\begin{equation}\label{eq:D-corr}
D^*=f\gamma\lambda^2\nu ,\qquad f=\lim_{\tau\to0}\frac{\langle X^2\rangle_\text{реал}}{\langle X^2\rangle_\text{случ}} .
\end{equation}
Грубая оценка: вакансия после обмена с атомом с вероятностью около
$1/z$ снова обменяется именно с ним, поэтому
$f\approx(1-1/z)^2\approx1-2/z$. Как видно из табл.~\ref{tab:corr-f}, эта
оценка неплоха.

\begin{table}[ht]
\centering
\caption{Корреляционный множитель $f$ для вакансионного механизма}
\label{tab:corr-f}
\small
\begin{tabular}{@{}lccc@{}}
\toprule
Решётка & $z$ & точный расчёт & $1-2/z$ \\
\midrule
двумерная квадратная & 4 & 0,467 & 0,50 \\
алмаз & 4 & 0,500 & 0,50 \\
простая кубическая & 6 & 0,653 & 0,67 \\
ОЦК & 8 & 0,727 & 0,75 \\
ГЦК, ГПУ & 12 & 0,781 & 0,83 \\
\bottomrule
\end{tabular}
\end{table}

Значение $f$ зависит от механизма:
\begin{itemize}
\item для механизмов, требующих дефекта по соседству (вакансионный,
непрямой межузельный), корреляция есть всегда;
\item прямой обмен и кольцевой механизм в чистом кубическом металле, а
также прямой межузельный механизм при малой концентрации внедрённых
атомов~— корреляции нет, $f\approx1$;
\item краудион блуждает одномерно и почти наверняка возвращается назад:
$f\approx0$;
\item дивакансионный механизм в плотной упаковке: $f\approx0{,}47$.
\end{itemize}
Для \emph{примесной} диффузии корреляция ещё важнее. Если вакансия
сильно связана с атомом примеси, она многократно меняется местами
именно с ним, и примесный атом топчется на месте, пока вакансия не
обменяется с атомом растворителя. В разбавленном ГЦК-растворе
приходится различать пять частот обмена вакансии: с атомом
растворителя вдали от примеси; с растворителем, соседним примеси (без
отрыва от неё, с отрывом, с подходом к ней) и с самим атомом примеси.
Множитель $f$ нельзя измерить прямо, но его можно найти из изотопного
эффекта.

\paragraph{Изотопный эффект и активационный объём}
Изотопы одного элемента с массами $m_\alpha$ и $m_\beta$ колеблются с
разными частотами ($\nu\propto m^{-1/2}$), поэтому диффундируют с
разной скоростью. Величина \term{изотопного эффекта}
\begin{equation}\label{eq:isotope}
E=\frac{D_\alpha^*/D_\beta^*-1}{(m_\beta/m_\alpha)^{1/2}-1}=f\,\Delta K ,
\end{equation}
где $0<\Delta K\le1$~— доля кинетической энергии, которую уносит
перескакивающий атом. Зная $f$, по $E$ находят $\Delta K$
(табл.~\ref{tab:isotope}); для примесей, наоборот, измеренное $E$
позволяет судить об $f$, а значит, о механизме.

Ещё один инструмент~— \emph{давление}. Из \eqref{eq:jump-freq}
$D^*\propto\exp(-G_f/kT)\exp(-G_m/kT)$, а $(\partial G/\partial p)_T=V$, поэтому
\begin{keyeq}\label{eq:act-volume}
\left(\frac{\partial\ln D^*}{\partial p}\right)_T=-\frac{\Delta V_f+\Delta V_m}{kT}\equiv-\frac{\Delta V^*}{kT},
\end{keyeq}
где $\Delta V^*$~— \term{активационный объём}: изменение объёма
кристалла при образовании дефекта ($\Delta V_f$) и при переходе атома в
седловую точку ($\Delta V_m$). При вакансионном механизме образование
вакансии увеличивает объём, $\Delta V^*>0$, и под давлением диффузия
\emph{замедляется}: сжатый образец «теряет» вакансии. Отрицательный
$\Delta V^*$~— указание на межузельный механизм. Эмпирически
$\Delta K\approx\Delta V^*/\Omega$ ($\Omega$~— атомный объём), и это
хорошо согласуется с опытом (табл.~\ref{tab:isotope}).

\begin{table}[ht]
\centering
\caption{Изотопный эффект и активационный объём самодиффузии}
\label{tab:isotope}
\small
\begin{tabular}{@{}llccccc@{}}
\toprule
Металл & Решётка & Изотопы & $E$ & $\Delta K=E/f$ & $\Delta V^*$, см$^3$/моль & $\Delta V^*/\Omega$ \\
\midrule
Ag & ГЦК & $^{105}$Ag, $^{111}$Ag & 0,67 & $0{,}86\pm0{,}05$ & 9,2 & 0,89 \\
$\gamma$-Fe & ГЦК & $^{55}$Fe, $^{59}$Fe & 0,58 & $0{,}74\pm0{,}08$ & --- & 0,77 \\
Na & ОЦК & $^{22}$Na, $^{24}$Na & 0,36 & $0{,}50\pm0{,}04$ & 12,4 & 0,52 \\
Au & ГЦК & & & & 7,2 & 0,70 \\
Pb & ГЦК & & & & 13,0 & 0,71 \\
\bottomrule
\end{tabular}
\end{table}

\paragraph{Пути ускоренной диффузии}
Диффузия~— процесс \emph{структурно чувствительный}: всё, что нарушает
идеальный порядок, её ускоряет. Области кристалла с нарушенным порядком
называют \term{путями ускоренной диффузии}.

\emph{Неравновесные вакансии.} Закалка, пластическая деформация,
фазовые превращения и облучение создают избыточные вакансии. Так как
$D^*=C_vD_v$ ($D_v$~— коэффициент диффузии вакансий), пересыщение
вакансиями ускоряет диффузию. Избыток уходит на стоки по закону
$\partial C_v/\partial t=(C_v^0-C_v)/\tau_v$ с временем релаксации
$\tau_v=AL^2/D_v$ ($L$~— расстояние до стока, $A$~— его эффективность).
Измеряя $\tau_v$ при разных температурах, находят энергию миграции
вакансий (п.~2.1). Если стоки~— только поверхность и границы зёрен,
избыточные вакансии заметно ускоряют диффузию.

\emph{Дислокации.} После фазовых превращений и деформации плотность
дислокаций на порядки выше, чем в отожжённом монокристалле. Вдоль
дислокации («трубки» диаметром в несколько векторов Бюргерса)
коэффициент диффузии на несколько порядков больше, чем в решётке, а
энергия активации $Q_d\approx(0{,}4$--$0{,}7)\,Q$. Чем ниже температура, тем
больше относительный вклад дислокаций. Если диффузионные зоны соседних
дислокаций перекрываются, работает приближение Харта:
$D_\text{эфф}=fD_d+(1-f)D_l$, где $f$~— доля времени, которую атом
проводит в трубках. Если не перекрываются~— анализ, аналогичный модели
Фишера для границ.

\emph{Границы зёрен.} Измельчение зерна ускоряет диффузию, и при
$T<(0{,}6$--$0{,}7)\Tpl$ график Аррениуса перестаёт быть прямой
(рис.~\ref{fig:gb-arrhenius}): на первый план выходит диффузия по
границам, у которой $D_g\gg D_l$, но $Q_g<Q_l$.

\begin{figure}[ht]
\centering
\begin{tikzpicture}
\begin{axis}[width=0.6\textwidth, height=5.2cm, xmin=1, xmax=2.2, ymin=-18.5, ymax=-5,
  xlabel={$\Tpl/T$}, ylabel={$\lg D$}, ytick=\empty, axis lines=left]
  \addplot[fmBlue, thick, domain=1:2.2] {-8-7.82*(x-1)};
  \addplot[fmRed, thick, domain=1:2.2] {-6-3.9*(x-1)};
  \addplot[black, very thick, dashed, domain=1:2.2, samples=60] {ln(10^(-8-7.82*(x-1))+10^(-10-3.9*(x-1)))/ln(10)};
  \node[font=\small, fmBlue, anchor=north east] at (axis cs:2.15,-16.6) {объём, $Q_l$};
  \node[font=\small, fmRed, anchor=south west] at (axis cs:1.5,-7.7) {граница, $Q_g$};
  \node[font=\small, anchor=south west] at (axis cs:1.62,-12.6) {поликристалл};
\end{axis}
\end{tikzpicture}
\caption{Температурная зависимость коэффициентов объёмной и
зернограничной диффузии (схема). Наклон прямой пропорционален энергии
активации. У поликристалла (пунктир) при высоких температурах
преобладает объёмная диффузия, при низких~— граничная}
\label{fig:gb-arrhenius}
\end{figure}

Разделить потоки по объёму и по границе позволяет \term{модель
Фишера} (рис.~\ref{fig:fisher}). Граница~— однородная пластина ширины
$\delta=2a$ между двумя зёрнами, перпендикулярная поверхности, с
$D_g\gg D_l$. Вещество быстро бежит вниз по границе и одновременно
«отсасывается» из неё в объём зёрен. Решение даёт
\[
C\propto\exp\!\left[-\frac{(4D_l/\pi t)^{1/4}}{(\delta D_g)^{1/2}}\,y\right]
\operatorname{erfc}\frac{x}{2\sqrt{D_lt}} .
\]
\begin{itemize}
\item Признак граничной диффузии: $\ln C$ линейно зависит от глубины
$y$, а не от $y^2$, как при объёмной, а наклон убывает как $t^{-1/4}$.
\item По наклону $\operatorname{tg}\alpha$ прямой $\ln C$~— $y$ находят
$\delta D_g/2=aD_g=(\operatorname{tg}^2\alpha)^{-1}(D_l/\pi t)^{1/2}$. Сам
по себе $D_g$ не определяется: ширину границы оценивают от 0,5~нм
(кристаллографическая ширина) до микрометров («диффузионная» ширина).
\item Итоги опытов: $D_g/D_l\approx10^3$--$10^5$, $Q_g/Q_l\approx0{,}35$--$0{,}7$.
Механизм граничной диффузии изучен хуже всего; обычно предполагают
вакансионный с $Q_g\approx E_m$ (С.\,Д.~Герцрикен): в границе вакансий
много, и энергию их образования можно не учитывать.
\end{itemize}

\begin{figure}[ht]
\centering
\begin{tikzpicture}[x=2.4cm, y=0.33cm, fisher/.style={fmBlue, line width=0.7pt}]
  \fill[fmRed!12] (-0.04,0) rectangle (0.04,-11.5);
  \draw[fmRed, line width=0.5pt] (-0.04,0) -- (-0.04,-11.5) (0.04,0) -- (0.04,-11.5);
  % сгенерировано tools/gen_diffusion.py — не править вручную
\draw[fisher] plot[smooth] coordinates {(-0.364,-0.000) (-0.363,-0.036) (-0.350,-0.098) (-0.336,-0.160) (-0.323,-0.220) (-0.309,-0.280) (-0.296,-0.340) (-0.282,-0.398) (-0.269,-0.456) (-0.255,-0.513) (-0.242,-0.569) (-0.229,-0.624) (-0.215,-0.679) (-0.202,-0.733) (-0.188,-0.786) (-0.175,-0.839) (-0.161,-0.891) (-0.148,-0.942) (-0.134,-0.992) (-0.121,-1.042) (-0.108,-1.091) (-0.094,-1.140) (-0.081,-1.187) (-0.067,-1.234) (-0.054,-1.281) (-0.040,-1.326) (-0.027,-1.371) (-0.013,-1.416) (-0.000,-1.460) (0.000,-1.460) (0.013,-1.416) (0.027,-1.371) (0.040,-1.326) (0.054,-1.281) (0.067,-1.234) (0.081,-1.187) (0.094,-1.140) (0.108,-1.091) (0.121,-1.042) (0.134,-0.992) (0.148,-0.942) (0.161,-0.891) (0.175,-0.839) (0.188,-0.786) (0.202,-0.733) (0.215,-0.679) (0.229,-0.624) (0.242,-0.569) (0.255,-0.513) (0.269,-0.456) (0.282,-0.398) (0.296,-0.340) (0.309,-0.280) (0.323,-0.220) (0.336,-0.160) (0.350,-0.098) (0.363,-0.036) (0.364,-0.000)};
\draw[fisher] plot[smooth] coordinates {(-0.660,-0.000) (-0.659,-0.012) (-0.645,-0.092) (-0.632,-0.170) (-0.618,-0.248) (-0.605,-0.325) (-0.592,-0.401) (-0.578,-0.477) (-0.565,-0.551) (-0.551,-0.625) (-0.538,-0.698) (-0.524,-0.771) (-0.511,-0.842) (-0.497,-0.913) (-0.484,-0.983) (-0.471,-1.052) (-0.457,-1.120) (-0.444,-1.187) (-0.430,-1.254) (-0.417,-1.320) (-0.403,-1.385) (-0.390,-1.450) (-0.376,-1.513) (-0.363,-1.576) (-0.350,-1.638) (-0.336,-1.700) (-0.323,-1.760) (-0.309,-1.820) (-0.296,-1.880) (-0.282,-1.938) (-0.269,-1.996) (-0.255,-2.053) (-0.242,-2.109) (-0.229,-2.164) (-0.215,-2.219) (-0.202,-2.273) (-0.188,-2.326) (-0.175,-2.379) (-0.161,-2.431) (-0.148,-2.482) (-0.134,-2.532) (-0.121,-2.582) (-0.108,-2.631) (-0.094,-2.680) (-0.081,-2.727) (-0.067,-2.774) (-0.054,-2.821) (-0.040,-2.866) (-0.027,-2.911) (-0.013,-2.956) (-0.000,-2.999) (0.000,-2.999) (0.013,-2.956) (0.027,-2.911) (0.040,-2.866) (0.054,-2.821) (0.067,-2.774) (0.081,-2.727) (0.094,-2.680) (0.108,-2.631) (0.121,-2.582) (0.134,-2.532) (0.148,-2.482) (0.161,-2.431) (0.175,-2.379) (0.188,-2.326) (0.202,-2.273) (0.215,-2.219) (0.229,-2.164) (0.242,-2.109) (0.255,-2.053) (0.269,-1.996) (0.282,-1.938) (0.296,-1.880) (0.309,-1.820) (0.323,-1.760) (0.336,-1.700) (0.350,-1.638) (0.363,-1.576) (0.376,-1.513) (0.390,-1.450) (0.403,-1.385) (0.417,-1.320) (0.430,-1.254) (0.444,-1.187) (0.457,-1.120) (0.471,-1.052) (0.484,-0.983) (0.497,-0.913) (0.511,-0.842) (0.524,-0.771) (0.538,-0.698) (0.551,-0.625) (0.565,-0.551) (0.578,-0.477) (0.592,-0.401) (0.605,-0.325) (0.618,-0.248) (0.632,-0.170) (0.645,-0.092) (0.659,-0.012) (0.660,-0.000)};
\draw[fisher] plot[smooth] coordinates {(-0.942,-0.000) (-0.941,-0.050) (-0.928,-0.147) (-0.914,-0.243) (-0.901,-0.339) (-0.887,-0.433) (-0.874,-0.527) (-0.861,-0.620) (-0.847,-0.712) (-0.834,-0.803) (-0.820,-0.894) (-0.807,-0.983) (-0.793,-1.072) (-0.780,-1.159) (-0.766,-1.246) (-0.753,-1.333) (-0.739,-1.418) (-0.726,-1.502) (-0.713,-1.586) (-0.699,-1.669) (-0.686,-1.751) (-0.672,-1.832) (-0.659,-1.912) (-0.645,-1.992) (-0.632,-2.070) (-0.618,-2.148) (-0.605,-2.225) (-0.592,-2.301) (-0.578,-2.377) (-0.565,-2.451) (-0.551,-2.525) (-0.538,-2.598) (-0.524,-2.670) (-0.511,-2.742) (-0.497,-2.813) (-0.484,-2.882) (-0.471,-2.952) (-0.457,-3.020) (-0.444,-3.087) (-0.430,-3.154) (-0.417,-3.220) (-0.403,-3.285) (-0.390,-3.350) (-0.376,-3.413) (-0.363,-3.476) (-0.350,-3.538) (-0.336,-3.600) (-0.323,-3.660) (-0.309,-3.720) (-0.296,-3.779) (-0.282,-3.838) (-0.269,-3.896) (-0.255,-3.952) (-0.242,-4.009) (-0.229,-4.064) (-0.215,-4.119) (-0.202,-4.173) (-0.188,-4.226) (-0.175,-4.279) (-0.161,-4.331) (-0.148,-4.382) (-0.134,-4.432) (-0.121,-4.482) (-0.108,-4.531) (-0.094,-4.580) (-0.081,-4.627) (-0.067,-4.674) (-0.054,-4.721) (-0.040,-4.766) (-0.027,-4.811) (-0.013,-4.856) (-0.000,-4.899) (0.000,-4.899) (0.013,-4.856) (0.027,-4.811) (0.040,-4.766) (0.054,-4.721) (0.067,-4.674) (0.081,-4.627) (0.094,-4.580) (0.108,-4.531) (0.121,-4.482) (0.134,-4.432) (0.148,-4.382) (0.161,-4.331) (0.175,-4.279) (0.188,-4.226) (0.202,-4.173) (0.215,-4.119) (0.229,-4.064) (0.242,-4.009) (0.255,-3.952) (0.269,-3.896) (0.282,-3.838) (0.296,-3.779) (0.309,-3.720) (0.323,-3.660) (0.336,-3.600) (0.350,-3.538) (0.363,-3.476) (0.376,-3.413) (0.390,-3.350) (0.403,-3.285) (0.417,-3.220) (0.430,-3.154) (0.444,-3.087) (0.457,-3.020) (0.471,-2.952) (0.484,-2.882) (0.497,-2.813) (0.511,-2.742) (0.524,-2.670) (0.538,-2.598) (0.551,-2.525) (0.565,-2.451) (0.578,-2.377) (0.592,-2.301) (0.605,-2.225) (0.618,-2.148) (0.632,-2.070) (0.645,-1.992) (0.659,-1.912) (0.672,-1.832) (0.686,-1.751) (0.699,-1.669) (0.713,-1.586) (0.726,-1.502) (0.739,-1.418) (0.753,-1.333) (0.766,-1.246) (0.780,-1.159) (0.793,-1.072) (0.807,-0.983) (0.820,-0.894) (0.834,-0.803) (0.847,-0.712) (0.861,-0.620) (0.874,-0.527) (0.887,-0.433) (0.901,-0.339) (0.914,-0.243) (0.928,-0.147) (0.941,-0.050) (0.942,-0.000)};
\draw[fisher] plot[smooth] coordinates {(-1.238,-0.000) (-1.237,-0.008) (-1.224,-0.125) (-1.210,-0.240) (-1.197,-0.355) (-1.183,-0.469) (-1.170,-0.582) (-1.156,-0.694) (-1.143,-0.805) (-1.129,-0.915) (-1.116,-1.025) (-1.103,-1.133) (-1.089,-1.241) (-1.076,-1.348) (-1.062,-1.454) (-1.049,-1.559) (-1.035,-1.663) (-1.022,-1.766) (-1.008,-1.868) (-0.995,-1.970) (-0.982,-2.070) (-0.968,-2.170) (-0.955,-2.269) (-0.941,-2.367) (-0.928,-2.464) (-0.914,-2.560) (-0.901,-2.656) (-0.887,-2.750) (-0.874,-2.844) (-0.861,-2.937) (-0.847,-3.029) (-0.834,-3.120) (-0.820,-3.211) (-0.807,-3.300) (-0.793,-3.389) (-0.780,-3.476) (-0.766,-3.563) (-0.753,-3.649) (-0.739,-3.735) (-0.726,-3.819) (-0.713,-3.903) (-0.699,-3.986) (-0.686,-4.068) (-0.672,-4.149) (-0.659,-4.229) (-0.645,-4.308) (-0.632,-4.387) (-0.618,-4.465) (-0.605,-4.542) (-0.592,-4.618) (-0.578,-4.694) (-0.565,-4.768) (-0.551,-4.842) (-0.538,-4.915) (-0.524,-4.987) (-0.511,-5.059) (-0.497,-5.130) (-0.484,-5.199) (-0.471,-5.268) (-0.457,-5.337) (-0.444,-5.404) (-0.430,-5.471) (-0.417,-5.537) (-0.403,-5.602) (-0.390,-5.667) (-0.376,-5.730) (-0.363,-5.793) (-0.350,-5.855) (-0.336,-5.917) (-0.323,-5.977) (-0.309,-6.037) (-0.296,-6.096) (-0.282,-6.155) (-0.269,-6.212) (-0.255,-6.269) (-0.242,-6.326) (-0.229,-6.381) (-0.215,-6.436) (-0.202,-6.490) (-0.188,-6.543) (-0.175,-6.596) (-0.161,-6.648) (-0.148,-6.699) (-0.134,-6.749) (-0.121,-6.799) (-0.108,-6.848) (-0.094,-6.897) (-0.081,-6.944) (-0.067,-6.991) (-0.054,-7.038) (-0.040,-7.083) (-0.027,-7.128) (-0.013,-7.173) (-0.000,-7.216) (0.000,-7.216) (0.013,-7.173) (0.027,-7.128) (0.040,-7.083) (0.054,-7.038) (0.067,-6.991) (0.081,-6.944) (0.094,-6.897) (0.108,-6.848) (0.121,-6.799) (0.134,-6.749) (0.148,-6.699) (0.161,-6.648) (0.175,-6.596) (0.188,-6.543) (0.202,-6.490) (0.215,-6.436) (0.229,-6.381) (0.242,-6.326) (0.255,-6.269) (0.269,-6.212) (0.282,-6.155) (0.296,-6.096) (0.309,-6.037) (0.323,-5.977) (0.336,-5.917) (0.350,-5.855) (0.363,-5.793) (0.376,-5.730) (0.390,-5.667) (0.403,-5.602) (0.417,-5.537) (0.430,-5.471) (0.444,-5.404) (0.457,-5.337) (0.471,-5.268) (0.484,-5.199) (0.497,-5.130) (0.511,-5.059) (0.524,-4.987) (0.538,-4.915) (0.551,-4.842) (0.565,-4.768) (0.578,-4.694) (0.592,-4.618) (0.605,-4.542) (0.618,-4.465) (0.632,-4.387) (0.645,-4.308) (0.659,-4.229) (0.672,-4.149) (0.686,-4.068) (0.699,-3.986) (0.713,-3.903) (0.726,-3.819) (0.739,-3.735) (0.753,-3.649) (0.766,-3.563) (0.780,-3.476) (0.793,-3.389) (0.807,-3.300) (0.820,-3.211) (0.834,-3.120) (0.847,-3.029) (0.861,-2.937) (0.874,-2.844) (0.887,-2.750) (0.901,-2.656) (0.914,-2.560) (0.928,-2.464) (0.941,-2.367) (0.955,-2.269) (0.968,-2.170) (0.982,-2.070) (0.995,-1.970) (1.008,-1.868) (1.022,-1.766) (1.035,-1.663) (1.049,-1.559) (1.062,-1.454) (1.076,-1.348) (1.089,-1.241) (1.103,-1.133) (1.116,-1.025) (1.129,-0.915) (1.143,-0.805) (1.156,-0.694) (1.170,-0.582) (1.183,-0.469) (1.197,-0.355) (1.210,-0.240) (1.224,-0.125) (1.237,-0.008) (1.238,-0.000)};
\draw[fisher] plot[smooth] coordinates {(-1.534,-0.000) (-1.533,-0.017) (-1.519,-0.154) (-1.506,-0.290) (-1.492,-0.424) (-1.479,-0.558) (-1.466,-0.691) (-1.452,-0.823) (-1.439,-0.954) (-1.425,-1.084) (-1.412,-1.214) (-1.398,-1.342) (-1.385,-1.469) (-1.371,-1.596) (-1.358,-1.721) (-1.345,-1.846) (-1.331,-1.970) (-1.318,-2.093) (-1.304,-2.214) (-1.291,-2.336) (-1.277,-2.456) (-1.264,-2.575) (-1.250,-2.693) (-1.237,-2.811) (-1.224,-2.927) (-1.210,-3.043) (-1.197,-3.158) (-1.183,-3.271) (-1.170,-3.384) (-1.156,-3.496) (-1.143,-3.608) (-1.129,-3.718) (-1.116,-3.827) (-1.103,-3.936) (-1.089,-4.043) (-1.076,-4.150) (-1.062,-4.256) (-1.049,-4.361) (-1.035,-4.465) (-1.022,-4.568) (-1.008,-4.671) (-0.995,-4.772) (-0.982,-4.873) (-0.968,-4.972) (-0.955,-5.071) (-0.941,-5.169) (-0.928,-5.267) (-0.914,-5.363) (-0.901,-5.458) (-0.887,-5.553) (-0.874,-5.647) (-0.861,-5.739) (-0.847,-5.831) (-0.834,-5.923) (-0.820,-6.013) (-0.807,-6.102) (-0.793,-6.191) (-0.780,-6.279) (-0.766,-6.366) (-0.753,-6.452) (-0.739,-6.537) (-0.726,-6.622) (-0.713,-6.705) (-0.699,-6.788) (-0.686,-6.870) (-0.672,-6.951) (-0.659,-7.031) (-0.645,-7.111) (-0.632,-7.189) (-0.618,-7.267) (-0.605,-7.344) (-0.592,-7.421) (-0.578,-7.496) (-0.565,-7.571) (-0.551,-7.645) (-0.538,-7.718) (-0.524,-7.790) (-0.511,-7.861) (-0.497,-7.932) (-0.484,-8.002) (-0.471,-8.071) (-0.457,-8.139) (-0.444,-8.207) (-0.430,-8.273) (-0.417,-8.339) (-0.403,-8.405) (-0.390,-8.469) (-0.376,-8.533) (-0.363,-8.596) (-0.350,-8.658) (-0.336,-8.719) (-0.323,-8.780) (-0.309,-8.840) (-0.296,-8.899) (-0.282,-8.957) (-0.269,-9.015) (-0.255,-9.072) (-0.242,-9.128) (-0.229,-9.183) (-0.215,-9.238) (-0.202,-9.292) (-0.188,-9.346) (-0.175,-9.398) (-0.161,-9.450) (-0.148,-9.501) (-0.134,-9.552) (-0.121,-9.601) (-0.108,-9.651) (-0.094,-9.699) (-0.081,-9.747) (-0.067,-9.794) (-0.054,-9.840) (-0.040,-9.886) (-0.027,-9.931) (-0.013,-9.975) (-0.000,-10.019) (0.000,-10.019) (0.013,-9.975) (0.027,-9.931) (0.040,-9.886) (0.054,-9.840) (0.067,-9.794) (0.081,-9.747) (0.094,-9.699) (0.108,-9.651) (0.121,-9.601) (0.134,-9.552) (0.148,-9.501) (0.161,-9.450) (0.175,-9.398) (0.188,-9.346) (0.202,-9.292) (0.215,-9.238) (0.229,-9.183) (0.242,-9.128) (0.255,-9.072) (0.269,-9.015) (0.282,-8.957) (0.296,-8.899) (0.309,-8.840) (0.323,-8.780) (0.336,-8.719) (0.350,-8.658) (0.363,-8.596) (0.376,-8.533) (0.390,-8.469) (0.403,-8.405) (0.417,-8.339) (0.430,-8.273) (0.444,-8.207) (0.457,-8.139) (0.471,-8.071) (0.484,-8.002) (0.497,-7.932) (0.511,-7.861) (0.524,-7.790) (0.538,-7.718) (0.551,-7.645) (0.565,-7.571) (0.578,-7.496) (0.592,-7.421) (0.605,-7.344) (0.618,-7.267) (0.632,-7.189) (0.645,-7.111) (0.659,-7.031) (0.672,-6.951) (0.686,-6.870) (0.699,-6.788) (0.713,-6.705) (0.726,-6.622) (0.739,-6.537) (0.753,-6.452) (0.766,-6.366) (0.780,-6.279) (0.793,-6.191) (0.807,-6.102) (0.820,-6.013) (0.834,-5.923) (0.847,-5.831) (0.861,-5.739) (0.874,-5.647) (0.887,-5.553) (0.901,-5.458) (0.914,-5.363) (0.928,-5.267) (0.941,-5.169) (0.955,-5.071) (0.968,-4.972) (0.982,-4.873) (0.995,-4.772) (1.008,-4.671) (1.022,-4.568) (1.035,-4.465) (1.049,-4.361) (1.062,-4.256) (1.076,-4.150) (1.089,-4.043) (1.103,-3.936) (1.116,-3.827) (1.129,-3.718) (1.143,-3.608) (1.156,-3.496) (1.170,-3.384) (1.183,-3.271) (1.197,-3.158) (1.210,-3.043) (1.224,-2.927) (1.237,-2.811) (1.250,-2.693) (1.264,-2.575) (1.277,-2.456) (1.291,-2.336) (1.304,-2.214) (1.318,-2.093) (1.331,-1.970) (1.345,-1.846) (1.358,-1.721) (1.371,-1.596) (1.385,-1.469) (1.398,-1.342) (1.412,-1.214) (1.425,-1.084) (1.439,-0.954) (1.452,-0.823) (1.466,-0.691) (1.479,-0.558) (1.492,-0.424) (1.506,-0.290) (1.519,-0.154) (1.533,-0.017) (1.534,-0.000)};

  \draw[line width=0.8pt] (-1.8,0) -- (1.8,0);
  \node[font=\small, anchor=south] at (0,0.2) {поверхность, $C=C_0$};
  \node[font=\small] at (-1.2,-7) {зерно};
  \node[font=\small] at (1.2,-7) {зерно};
  \draw[-{Stealth[length=2mm]}] (1.95,0) -- (1.95,-4) node[right, font=\small]{$y$};
  \draw[-{Stealth[length=2mm]}] (0.55,-10.5) -- (1.2,-10.5) node[right, font=\small]{$x$};
  \node[font=\scriptsize, fmRed, anchor=west] at (0.06,-11) {граница ($D_g$)};
\end{tikzpicture}
\caption{Модель Фишера: линии равной концентрации (расчёт). Вдоль
границы вещество проникает гораздо глубже, чем в объём зёрен, и
изоконцентраты вытягиваются вдоль неё}
\label{fig:fisher}
\end{figure}

\emph{Поверхность.} Поверхность кристалла состоит из атомно-гладких
\emph{террас}, разделённых \emph{ступеньками} с \emph{изломами}. На ней
есть адсорбированные атомы (адатомы) и адвакансии, концентрация
которых $n_a=n_0\exp(-E_{af}/kT)$. Коэффициент поверхностной диффузии,
как и объёмной,~— произведение концентрации подвижных дефектов на их
подвижность:
$D_s=n_0P_a\nu_a\lambda_a^2\exp[-(E_{af}+E_{ma})/kT]$. Энергия миграции
адатома мала (на плоскости (111) ГЦК-кристалла около $1/20$ теплоты
сублимации). Адатомы буквально катятся по поверхности, смещаясь за
один акт на много межатомных расстояний,~— механизм «перекати-поле»
(\emph{rolling stone}). Индивидуальные скачки преобладают лишь при
$T<0{,}15\Tpl$; при более высоких температурах работают коллективные
механизмы. Поверхностная диффузия сглаживает рельеф (движущая сила~—
уменьшение поверхностной энергии), но там, где на поверхность выходит
граница зерна или дислокация, наоборот, развиваются канавки и ямки
термического травления.

\begin{glavnoe}
Иерархия скоростей диффузии: поверхность $>$ границы зёрен $>$ дислокации
$>$ объём. Энергии активации идут в обратном порядке, поэтому чем ниже
температура, тем больше роль «быстрых путей». Именно поэтому
мелкозернистые и деформированные материалы при умеренных температурах
ведут себя совсем иначе, чем монокристаллы.
\end{glavnoe}

\paragraph{«Нормальная» и «аномальная» самодиффузия}
К 1950-м годам сложилось представление о \emph{нормальной}
самодиффузии металлов: уравнение Аррениуса с $D_0\approx10^{-2}$--$10$~см$^2$/с
и энергией активации, пропорциональной температуре плавления,
$Q\approx18\,R\Tpl$\footnote{В оригинале это правило записано как
$Q=(35{,}3\pm1{,}5)\Tpl$ для ОЦК- и $(36{,}5\pm1{,}5)\Tpl$ для
плотноупакованных металлов без указания единиц; коэффициенты
соответствуют $Q$ в кал/моль при $\Tpl$ в кельвинах, т.\,е.
$Q/R\Tpl\approx18$. Аналогично ниже $(18{,}5\pm2{,}3)\Tpl$~кал/моль
$\approx9\,R\Tpl$.}, вакансионным механизмом, $f$ из
табл.~\ref{tab:corr-f} и активационным объёмом $(0{,}5$--$0{,}9)\,\Omega$.

Затем выяснилось, что ОЦК-фазы $\beta$-Zr, $\beta$-Ti, $\beta$-Hf,
$\gamma$-U, $\varepsilon$-Pu и некоторых лантаноидов~— важнейших
реакторных металлов~— ведут себя \emph{аномально}
(рис.~\ref{fig:anomalous}):
\begin{itemize}
\item малые параметры: $D_0\approx10^{-5}$--$10^{-2}$~см$^2$/с,
$Q\approx9\,R\Tpl$;
\item у точки плавления $D$ на порядок больше, чем у нормальных металлов;
\item сильно искривлённый график Аррениуса (у $\beta$-Zr и $\beta$-Ti
часть искривления устраняется правильным выбором времени отжига и
отжигом неравновесных дефектов);
\item у $\varepsilon$-Pu и $\delta$-Ce активационный объём отрицателен, а
изотопный эффект у $\beta$-Zr и $\beta$-Ti с понижением температуры
падает до $E\le0{,}3$.
\end{itemize}
Эти факты объясняют участием нескольких механизмов: кроме обычных
вакансионных NN-скачков~— межузельные гантели $\langle110\rangle$,
дивакансии, NNN-скачки. Межузельный вклад правдоподобен: у переходных
металлов с $d^1$- и $d^2$-оболочками аномально низкие модули
упругости, поэтому энергия образования межузельного атома невелика.
Кроме того, ускорение связывают с быстро диффундирующими комплексами
«вакансия~— примесный атом» (особенно с малыми атомами примесей
группы железа); тогда
\begin{equation}\label{eq:Q-complex}
Q=E_f+E_m-E_b ,
\end{equation}
где $E_b$~— энергия связи вакансии с примесью.

\begin{figure}[ht]
\centering
\begin{tikzpicture}
\begin{axis}[width=0.6\textwidth, height=5.2cm, xmin=1, xmax=1.8, ymin=-16, ymax=-6,
  xlabel={$\Tpl/T$}, ylabel={$\lg D$, см$^2$/с}, axis lines=left]
  \addplot[fmBlue, thick, domain=1:1.8] {(-18*x)/ln(10)};
  \addplot[fmRed, thick, domain=1:1.8] {-3+(-9.3*x)/ln(10)};
  \node[font=\small, fmBlue, anchor=north east] at (axis cs:1.78,-13.6) {нормальные};
  \node[font=\small, fmRed, anchor=south west] at (axis cs:1.3,-8.5) {аномальные (ОЦК Zr, Ti, U, Pu)};
\end{axis}
\end{tikzpicture}
\caption{Самодиффузия «нормальных» ($D_0\approx1$~см$^2$/с,
$Q\approx18R\Tpl$) и «аномальных» ($D_0\approx10^{-3}$~см$^2$/с,
$Q\approx9R\Tpl$) металлов в приведённых координатах (оценка по
эмпирическим правилам). У аномальных металлов диффузия быстрее и слабее
зависит от температуры}
\label{fig:anomalous}
\end{figure}
